% Part: incompleteness
% Chapter: arithmetization-syntax
% Section: proofs-in-nd

\documentclass[../../../include/open-logic-section]{subfiles}

\begin{document}

\olfileid{inc}{art}{pnd}
\olsection{नैसर्गिक निगमनातील \usetoken{P}{derivation}}

\begin{explain}
!!{derivation}s चे अंकगणितीकरण करण्यासाठी !!{derivation}s संख्यांनी
दर्शवावे लागतात. प्रत्येक निगमनाला एक किंवा दोन लेबले असलेल्या
!!{formula}s चे वृक्ष म्हणजे !!{derivation}s असल्यामुळे, पुनरावर्ती
निरूपण हा सहज मार्ग आहे. आपण !!{derivation} एका क्रमित समूहाने दर्शवतो.
त्याचे घटक म्हणजे अखेरच्या निगमनाच्या आधारविधानांकडे नेणाऱ्या तात्काळ
उप-!!{derivation}s ची संख्या, त्या उप-!!{derivation}s ची निरूपणे,
अंतिम !!{formula}, अखेरच्या निगमनाचे मुक्तता-लेबल आणि त्या निगमनाचा
प्रकार दर्शवणारी संख्या.
\end{explain}

\begin{defn}
  $\delta$ ही नैसर्गिक निगमनातील !!a{derivation} असेल, तर
  $\Gn{\delta}$ ची विगमनाने व्याख्या पुढीलप्रमाणे करतो:
  \begin{enumerate}
  \item
    $\delta$ मध्ये फक्त गृहीतक~$!A$ असेल, तर $\Gn{\delta}$ म्हणजे
    $\tuple{0, \Gn{!A}, n}$. हे गृहीतक !!{undischarged} असेल,
    तर~$n$ ही संख्या~$0$ असते; अन्यथा ती त्याचे संख्यात्मक लेबल असते.
  \item
    $\delta$ चा शेवट शून्य, एक, दोन किंवा तीन आधारविधानांच्या
    निगमनाने होत असेल, तर $\Gn{\delta}$ अनुक्रमे
    \begin{align*}
      & \tuple{0, \Gn{!A}, n, k}, \\
      & \tuple{1, \Gn{\delta_1}, \Gn{!A}, n, k}, \\
      & \tuple{2, \Gn{\delta_1}, \Gn{\delta_2}, \Gn{!A}, n, k}, \text{ किंवा}\\
      & \tuple{3, \Gn{\delta_1}, \Gn{\delta_2}, \Gn{\delta_3}, \Gn{!A}, n,
        k},
    \end{align*}
    असते. येथे $\delta_1$, $\delta_2$, $\delta_3$ या $\delta$ मधील
    अखेरच्या निगमनाच्या आधारविधानांवर संपणाऱ्या उप-!!{derivation}s आहेत;
    $!A$ हा $\delta$ मधील त्या निगमनाचा निष्कर्ष आहे; $n$ हे त्या निगमनाचे
    मुक्तता-लेबल आहे (निगमनाने एकही गृहीतक मुक्त होत नसेल, तर~$0$);
    आणि अखेरच्या निगमनात वापरलेल्या नियमानुसार $k$ पुढील तक्त्यात
    दिले आहे.

    \begin{tabular}{lccccccc}
      \text{नियम:} & \Intro{\land} & \Elim{\land} & \Intro{\lor} & \Elim{\lor} \\
      $k$: & 1 & 2 & 3 & 4  \\[1.5ex]
      \text{नियम:} &  \Intro{\lif} & \Elim{\lif} & \Intro{\lnot} & \Elim{\lnot} \\
      $k$: & 5 & 6 & 7 & 8 \\[1.5ex]
      \text{नियम:}  &  \FalseInt & \FalseCl & \Intro{\lforall} & \Elim{\lforall} \\
      $k$: & 9 & 10 & 11 & 12 \\[1.5ex]
      \text{नियम:} & \Intro{\lexists} & \Elim{\lexists} & \Intro{\eq} & \Elim{\eq} \\
      $k$: & 13 & 14 & 15 & 16
    \end{tabular}
  \end{enumerate}
\end{defn}

\begin{ex}
  पुढील अगदी सोपी !!{derivation} पाहा:
  \begin{prooftree}
    \AxiomC{$\Discharge{!A \land !B}{1}$}
    \RightLabel{\Elim{\land}}
    \UnaryInfC{$!A$}
    \DischargeRule{\Intro{\lif}}{1}
    \UnaryInfC{$(!A \land !B) \lif !A$}
  \end{prooftree}
  गृहीतकाची गोडेल संख्या $d_0 = \tuple{0,
    \Gn{(!A \land !B)}, 1}$ असेल. $\Elim{\land}$ च्या निष्कर्षावर
  संपणाऱ्या !!{derivation} ची गोडेल संख्या $d_1 = \tuple{1,
     d_0, \Gn{!A}, 0, 2}$ असेल. येथे $1$ हे $\Elim{\land}$ ला एकच
  आधारविधान असल्यामुळे आले आहे; निष्कर्ष~$!A$ ची गोडेल संख्या पुढे
  आहे; $0$ हे एकही गृहीतक मुक्त होत नसल्यामुळे आले आहे; आणि $2$
  हा $\Elim{\land}$ चा संकेतांक आहे. मग संपूर्ण !!{derivation} ची
  गोडेल संख्या $\tuple{1, d_1, \Gn{((!A \land !B) \lif
      !A)}, 1, 5}$, म्हणजेच
  \[
  \tuple{1, \tuple{1, \tuple{0, \Gn{(!A \land !B)}, 1}, \Gn{!A}, 0, 2},
    \Gn{((!A \land !B) \lif !A)}, 1, 5}.
  \]
\end{ex}

\begin{explain}
!!{derivation}s चे निरूपण ठरल्यानंतर, अशा !!{derivation}s च्या
गोडेल संख्यांवर आदिम पुनरावर्ती रीतीने क्रिया करता येतात आणि त्यांचे
आवश्यक गुणधर्म व संबंध व्यक्त करता येतात, हेही दाखवावे लागेल.
काही क्रिया सोप्या आहेत. उदाहरणार्थ, !!a{derivation} ची गोडेल
संख्या~$d$ दिली असता, $\fn{EndFmla}(d) = (d)_{(d)_0+1}$ हे
तिच्या अंतिम !!{formula} ची गोडेल संख्या देते;
$\fn{DischargeLabel}(d) = (d)_{(d)_0 + 2}$ हे मुक्तता-लेबल देते;
आणि $\fn{LastRule}(d) = (d)_{(d)_0 + 3}$ हे अखेरच्या निगमनाचा
प्रकार दर्शवणारी संख्या देते. काही इतर क्रिया अधिक कठीण आहेत;
त्या कशा कराव्यात याची किमान रूपरेषा आपण पाहू. $\delta$ ही
$\Gamma$ पासून $!A$ ची !!a{derivation} आहे, हा संबंध
$\delta$ आणि $!A$ यांच्या गोडेल संख्यांचा आदिम पुनरावर्ती संबंध
आहे, हे दाखवणे आपले उद्दिष्ट आहे.
\end{explain}

\begin{prop}
  पुढील संबंध आदिम पुनरावर्ती आहेत:
  \begin{enumerate}
  \item $!A$ हे $\delta$ मध्ये लेबल~$n$ असलेले गृहीतक म्हणून येते.
  \item $\delta$ मधील लेबल~$n$ असलेली सर्व गृहीतके $!A$ या रूपाची
    असतात (म्हणजे $\delta$ मध्ये लेबल~$n$ वापरून गृहीतक~$!A$
    !!{discharge} करता येते).
  \end{enumerate}
\end{prop}

\begin{proof}
  !!{formula}s च्या गोडेल संख्या आणि !!{derivation}s च्या गोडेल
  संख्या यांच्यातील अनुरूप संबंध आदिम पुनरावर्ती आहेत, हे दाखवायचे आहे.
  \begin{enumerate}
  \item $\fn{Assum}(x, d, n)$ आदिम पुनरावर्ती आहे, हे दाखवू.
    गोडेल संख्या~$d$ असलेल्या !!{derivation} मधील लेबल~$n$ च्या
    गृहीतकाची गोडेल संख्या~$x$ असेल, तर हा संबंध खरा असतो.
    गोडेल संख्या~$\tuple{0, x, n}$ असलेली !!{derivation} ही
    $d$ ची उप-!!{derivation} असेल, तर आणि तरच असे घडते. येथे
    !!{derivation}s चे आपले सांकेतीकरण हे
    \olref[cmp][rec][tre]{sec} मध्ये दिलेल्या वृक्षांच्या
    सांकेतीकरणाचे विशेष प्रकरण आहे, हे लक्षात घ्या. म्हणून आदिम
    पुनरावर्ती फलन $\fn{SubtreeSeq}(d)$ हे $d$ च्या सर्व
    उप-!!{derivation}s च्या गोडेल संख्यांची क्रमिका देते
    ($d$ च्या या क्रमिकेत पुनरुक्तीही असू शकतात). म्हणून व्याख्या अशी करता येते:
    \[
    \fn{Assum}(x, d, n) \defiff \bexists{i<\len{\fn{SubtreeSeq}(d)}}{(\fn{SubtreeSeq}(d))_i =
      \tuple{0, x, n}}.
    \]
  \item $\fn{Discharge}(x, d, n)$ आदिम पुनरावर्ती आहे, हे दाखवू.
    गोडेल संख्या~$d$ असलेल्या !!{derivation} मधील लेबल~$n$
    असलेली सर्व गृहीतके गोडेल संख्या~$x$ असलेल्या !!{formula}
    चीच असतील, तर हा संबंध खरा असतो. तो
    $\bforall{y<d}{(\fn{Assum}(y, d, n) \lif y
      = x)}$ असेल, तर आणि तरच खरा असतो.
  \end{enumerate}
\end{proof}

\begin{prop}
  \ollabel{prop:followsby}
  गोडेल संख्या~$d$ असलेल्या !!{derivation}~$\delta$ मधील अखेरचे
  निगमन योग्य असेल, तर आणि तरच खरा असणारा $\fn{Correct}(d)$ गुणधर्म
  आदिम पुनरावर्ती आहे.
\end{prop}

\begin{proof}
  प्रत्येक निगमन नियम~$R$ साठी $\fn{FollowsBy}_R(d)$ हा संबंध
  आदिम पुनरावर्ती आहे, हे दाखवावे लागेल. $d$ ही !!{derivation}~$\delta$
  ची गोडेल संख्या असेल आणि $\delta$ चे अंतिम !!{formula} $\delta$ च्या
  तात्काळ उप-!!{derivation}s पासून $R$ च्या योग्य वापराने
  निष्पन्न होत असेल, तर आणि तरच $\fn{FollowsBy}_R(d)$ खरे असते.

  \Intro{\land} नियमाचे उदाहरण सोपे आहे. $\delta$ चा शेवट
  $\Intro{\land}$ च्या योग्य वापराने होत असेल, तर तिचे रूप असे असते:
  \begin{prooftree}
    \AxiomC{}
    \RightLabel{$\delta_1$}
    \DeduceC{$!A$}

    \AxiomC{}
    \RightLabel{$\delta_2$}
    \DeduceC{$!B$}

    \RightLabel{\Intro\land}
    \BinaryInfC{$!A \land !B$}
  \end{prooftree}
  मग $\delta$ ची गोडेल संख्या~$d$ ही
  $\tuple{2, d_1, d_2, \Gn{(!A \land !B)}, 0, k}$ असते. येथे
  $\fn{EndFmla}(d_1) = \Gn{!A}$,
  $\fn{EndFmla}(d_2) = \Gn{!B}$, $n=0$ आणि $k=1$.
  म्हणून $\fn{FollowsBy}_{\Intro\land}(d)$ ची व्याख्या अशी करता येते:
  \begin{multline*}
    (d)_0 = 2 \land \fn{DischargeLabel}(d) = 0 \land \fn{LastRule}(d) = 1 \land {}\\
    \fn{EndFmla}(d) = {}\\
    \Gn{(} \concat \fn{EndFmla}((d)_1) \concat \Gn{\land}
    \concat \fn{EndFmla}((d)_2) \concat \Gn{)}.
  \end{multline*}

  आणखी एक सोपे उदाहरण $\Intro\eq$ नियमाचे आहे. त्याला
  आधारविधाने नसतात, म्हणून गृहीतकांप्रमाणेच $(d)_0 = 0$ असते.
  त्याचे मुक्तता-लेबलही नसते, म्हणजे $n=0$. मात्र $!A$ हे
  एखाद्या बंद पद~$t$ साठी $\eq[t][t]$ या रूपाचे असले पाहिजे.
  येथे आदिम पुनरावर्ती व्याख्या अशी आहे:
  \begin{multline*}
    (d)_0 = 0 \land \fn{DischargeLabel}(d) = 0 \land {}\\
    \bexists{t<d}{(\fn{ClTerm}(t) \land
      \fn{EndFmla}(d) = {}\\
      \Gn{{\eq}(} \concat t \concat \Gn{,} \concat t \concat \Gn{)})}.
  \end{multline*}

  अधिक गुंतागुंतीचे उदाहरण म्हणून $\fn{FollowsBy}_{\Intro{\lif}}(d)$
  पाहू. $\delta$ चे अंतिम !!{formula} $(!A \lif !B)$ या रूपाचे,
  $\delta_1$ चे अंतिम !!{formula}~$!B$ असेल आणि $\delta$ मधील
  लेबल~$n$ असलेले प्रत्येक गृहीतक $!A$ या रूपाचे असेल,
  तर आणि तरच हा संबंध खरा असतो. हे आदिम पुनरावर्ती रीतीने
  पुढीलप्रमाणे व्यक्त करता येते:
  \begin{multline*}
    (d)_0 = 1 \land {}\\
    \bexists{a<d}{(\fn{Discharge}(a, (d)_1, \fn{DischargeLabel}(d)) \land {}}\\
      \fn{EndFmla}(d) = (\Gn{(} \concat a \concat \Gn{\lif}
      \concat \fn{EndFmla}((d)_1) \concat \Gn{)}))
  \end{multline*}
  (येथे $a$ ही $!A$ ची गोडेल संख्या समजा.)

  आणखी एका उदाहरणासाठी \Intro{\lexists} पाहू. $\delta$ मधील
  अखेरचे निगमन योग्य असेल, तर आणि तरच !!a{formula}~$!A$,
  बंद पद~$t$ आणि !!a{variable}~$x$ अशी असतात की
  $\Subst{!A}{t}{x}$ हे !!{derivation}~$\delta_1$ चे अंतिम
  !!{formula} असते आणि $\lexists[x][!A]$ हा अखेरच्या निगमनाचा
  निष्कर्ष असतो. म्हणून $\fn{FollowsBy}_{\Intro{\lexists}}(d)$
  पुढील अट पूर्ण झाली तर आणि तरच खरे असते:
  \begin{multline*}
    (d)_0 = 1 \land \fn{DischargeLabel}(d) = 0 \land {} \\
    \bexists{a < d}{\bexists{x<d}{\bexists{t<d}{
          (\fn{ClTerm}(t) \land \fn{Var}(x)  \land {}}}}\\
    \fn{Subst}(a,t,x) = \fn{EndFmla}((d)_1) \land
    \fn{EndFmla}(d) = (\Gn{\lexists} \concat x \concat a)).
  \end{multline*}

नेमका गाठीचा आकार तपासण्यासाठी मागील विभागातील वैध क्रमिकासंकेतांकाची
अट वापरा. गृहीतकाला तीन घटक असतात; नियमगाठीला आधारविधानांच्या
संख्येपेक्षा चार अधिक घटक असतात:
\[
\fn{NDNode}(d) \defiff \fn{Seq}(d) \land
[((d)_0=0 \land \len{d}=3) \lor ((d)_0\le3 \land \len{d}=(d)_0+4)].
\]
  मग $\fn{Correct}(d)$ ची व्याख्या अशी करतो:
  \begin{multline*}
    \fn{NDNode}(d) \land \fn{Sent}(\fn{EndFmla}(d)) \land {}\\
    [(\fn{LastRule}(d) = 1 \land
    \fn{FollowsBy}_{\Intro\land}(d)) \lor \dots \lor {}\\
    (\fn{LastRule}(d) = 16 \land \fn{FollowsBy}_{\Elim\eq}(d)) \lor {}\\
    \bexists{n<d}{\bexists{x<d}{(d = \tuple{0, x, n})}}].
  \end{multline*}
  पहिली ओळ संकेतांकाचा नेमका गाठीचा आकार आणि $d$ चे अंतिम !!{formula} वाक्य आहे, याची खात्री
  करते. शेवटची ओळ $d$ मध्ये फक्त गृहीतक असण्याचा प्रसंग हाताळते.
\end{proof}

\begin{prob}
  \olref[inc][art][pnd]{prop:followsby} प्रमाणे पुढील गुणधर्म
  परिभाषित करा:
  \begin{enumerate}
  \item $\fn{FollowsBy}_{\Elim{\lif}}(d)$,
  \item $\fn{FollowsBy}_{\Elim{\eq}}(d)$,
  \item $\fn{FollowsBy}_{\Elim{\lor}}(d)$,
  \item $\fn{FollowsBy}_{\Intro{\lforall}}(d)$.
  \end{enumerate}
  शेवटच्या गुणधर्मासाठी, गोडेल संख्या~$d$ असलेल्या
  !!{derivation} मधील अखेरचे निगमन आयगेन चराची अट पूर्ण करते
  का, हे आदिम पुनरावर्ती रीतीने तपासता येते, हेही दाखवा.
  म्हणजेच $\Intro{\lforall}$ निगमनाचा आयगेन चर~$a$ हा
  $d$ च्या अंतिम !!{formula} मध्ये किंवा $d$ च्या मुक्त न केलेल्या
  गृहीतकात येत नाही. यासाठी
  \olref[inc][art][pnd]{prop:openassum} मधील आदिम पुनरावर्ती
  विधेय $\fn{OpenAssum}$ वापरू शकता.
\end{prob}

\begin{prop}
  \ollabel{prop:deriv}
  $d$ ही योग्य !!{derivation}~$\delta$ ची गोडेल संख्या असेल,
  तर खरा असणारा $\fn{Deriv}(d)$ संबंध आदिम पुनरावर्ती आहे.
\end{prop}

\begin{proof}
  !!^a{derivation}~$\delta$ मधील प्रत्येक निगमन एखाद्या नियमाचा
  योग्य वापर असेल, म्हणजेच तिच्या प्रत्येक उप-!!{derivation}s चा
  शेवट योग्य निगमनाने होत असेल, तर ती योग्य असते. म्हणून
  $\fn{Deriv}(d)$ तर आणि तरच
  \[
  \bforall{i<\len{\fn{SubtreeSeq}(d)}}{\fn{Correct}((\fn{SubtreeSeq}(d))_i)}
  \]
\end{proof}

\begin{prop}
  \ollabel{prop:openassum} $z$ ही गोडेल संख्या
  !!a{undischarged} गृहीतक~$!A$ ची असेल, आणि हे गृहीतक
  गोडेल संख्या~$d$ असलेल्या !!{derivation}~$\delta$ मध्ये असेल,
  तर खरा असणारा $\fn{OpenAssum}(z, d)$ संबंध आदिम पुनरावर्ती आहे.
\end{prop}

\begin{proof}
  गृहीतकाच्या एखाद्या आस्थितीला शून्येतर लेबल~$n$ असेल आणि ती
  $\delta$ च्या अशा उप-!!{derivation} मध्ये येत असेल जिचा शेवट
  मुक्तता-लेबल~$n$ असलेल्या नियमाने होतो, तर ती आस्थिती
  !!{discharged} असते. म्हणून अशा गृहीतकाची किमान एक आस्थिती
  $\delta$ मध्ये !!{discharged} नसेल, तर $!A$ हे
  $\delta$ चे !!a{undischarged} गृहीतक असते. येथे सावध राहिले
  पाहिजे: $\delta$ मध्ये $!A$ च्या !!{discharged} आणि
  !!{undischarged} अशा दोन्ही आस्थिती असू शकतात.

  $\delta_0$, \dots, $\delta_k$ ही क्रमिका पाहा. येथे $\delta_0 =
  \delta$, $\delta_k$ हे गृहीतक $\Discharge{!A}{n}$ आहे
  (एखाद्या~$n$ साठी), आणि $\delta_{i+1}$ ही $\delta_i$ ची तात्काळ
  उप-!!{derivation} आहे. लेबल शून्य असेल, तर हे गृहीतक व्याख्येनेच
  मुक्त न केलेले असते. अन्यथा, अशा क्रमिकेत गृहीतकाच्या आधीच्या कोणत्याही
  $\delta_i$ चा शेवट मुक्तता-लेबल~$n$ असलेल्या निगमनाने होत नसेल,
  तर $!A$ हे $\delta$ चे !!a{undischarged} गृहीतक आहे.

  आदिम पुनरावर्ती फलन $\fn{SubtreeSeq}(d)$ हे $\delta$ च्या
  सर्व उप-!!{derivation}s च्या गोडेल संख्यांची क्रमिका देते.
  $\delta$ च्या मूळापासून गृहीतकापर्यंतचा उप-!!{derivation}s चा मार्ग त्यातून
  निवडता येतो. अशी उपक्रमिका ओळखणे हा आदिम पुनरावर्ती संबंध आहे:
  $\fn{Subseq}(s, s')$ तर आणि तरच
  $\bforall{i<\len{s}}{\lexists[j<\len{s'}][(s)_i = (s')_j]}$.
  येथे दिलेली अट प्रत्येक घटकाचे सदस्यत्व तपासते; ती स्वतःहून
  क्रम-संरक्षण व्यक्त करत नाही. तात्काळ उप-!!{derivation} असण्याचा
  संबंधही आदिम पुनरावर्ती आहे: $\fn{Subderiv}(d, d')$ तर आणि तरच
  $\bexists{j<(d')_0}{d = (d')_{j+1}}$. म्हणून
  $\fn{OpenAssum}(z, d)$ ची व्याख्या अशी करता येते:
  \begin{multline*}
    \bexists{s<1+\fn{sequenceBound}(d,d)}{(\fn{Seq}(s) \land \len{s}>0 \land {}\\
      \fn{Subseq}(s, \fn{SubtreeSeq}(d))
      \land (s)_0 = d \land {}} \\
    \bexists{n<d}{((s)_{\len{s} \tsub 1} = \tuple{0, z, n} \land {}}\\
      \bforall{i<(\len{s} \tsub 1)}{(\fn{Subderiv}((s)_{i+1}, (s)_i) \land {}}\\
      (n=0 \lor \fn{DischargeLabel}((s)_i) \neq n)))).
  \end{multline*}
मार्गात प्रत्येक पुढचा गाठीचा संकेतांक आधीच्यापेक्षा लहान असतो.
म्हणून वैध मार्गात जास्तीत जास्त $d$ गाठी आणि प्रत्येक संकेतांक
जास्तीत जास्त $d$ असतो; वरची आदिम पुनरावर्ती मर्यादा पुरेशी आहे.
%READERNOTE{SOL6-A066 — मूळ OpenAssum च्या मार्गसंकेतांकासाठी पूर्ण subtree क्रमिकेचा code हा strict upper bound घेतला आहे. फक्त एक गृहीतक असलेल्या निष्पत्तीसाठी आवश्यक singleton मार्गाचा code त्या क्रमिकेच्या code इतकाच असतो; म्हणून मूळ कसोटी त्या गृहीतकाला उघडे मानत नाही. येथे वैध nonempty मार्गाला सर्वसाधारण sequence bound दिला. Exact node guard ने extra fields वगळले, आणि Assum ने प्रत्यक्ष subtree क्रमिकेच्या पूर्ण लांबीपर्यंत शोध घेतला.}
\end{proof}

\begin{prop}
  \ollabel{prop:prf-prim-rec}
  समजा $\Gamma$ हा !!{sentence}s चा आदिम पुनरावर्ती संच आहे.
  मग $\Prf[\Gamma](x, y)$ हा संबंध आदिम पुनरावर्ती आहे. तो असे व्यक्त
  करतो: $x$ हा $\Gamma$ मधील !!{undischarged} गृहीतकांपासून
  $!A$ च्या !!a{derivation}~$\delta$ चा संकेतांक आहे, आणि
  $y$ ही $!A$ ची गोडेल संख्या आहे.
\end{prop}

\begin{proof}
  समजा ``$y \in \Gamma$'' हे $R_\Gamma(y)$ या आदिम पुनरावर्ती
  विधेयाने दिले आहे. $\Prf[\Gamma](x, y)$ आदिम पुनरावर्ती आहे,
  हे दाखवायचे आहे. $y$ ही वाक्य~$!A$ ची गोडेल संख्या असेल,
  आणि $x$ हा अंतिम !!{formula}~$!A$ असलेल्या नैसर्गिक निगमनातील
  !!{derivation} चा संकेतांक असेल, तसेच तिची सर्व
  !!{undischarged} गृहीतके $\Gamma$ मध्ये असतील, तर आणि तरच हा
  संबंध खरा असतो.

  \olref{prop:deriv} नुसार, $x$ ही नैसर्गिक निगमनातील योग्य
  !!{derivation}~$\delta$ ची गोडेल संख्या असेल, तर आणि तरच खरा
  असणारा $\fn{Deriv}(x)$ गुणधर्म आदिम पुनरावर्ती आहे. त्यामुळे
  $\Prf[\Gamma](x, y)$ ची व्याख्या अशी करता येते:
  \begin{align*}
    \Prf[\Gamma](x, y) \defiff {}
    & \fn{Deriv}(x) \land \fn{EndFmla}(x) = y \land {} \\
    & \bforall{z < x}{(\fn{OpenAssum}(z, x) \lif R_\Gamma(z))}.
  \end{align*}
\end{proof}

\end{document}
